Sea $u(x, y)$ una función con derivadas parciales de segundo orden continuas que satisface la ecuación de Laplace $u_{xx} + u_{yy} = 0$. Usando el cambio de variables $x = r\cos\theta$, $y = r\sin\theta$ y la regla de la cadena, demuestre que $u$ satisface la ecuación de Laplace en coordenadas polares:
\[
\frac{\partial^2 u}{\partial r^2} + \frac{1}{r}\frac{\partial u}{\partial r} + \frac{1}{r^2}\frac{\partial^2 u}{\partial \theta^2} = 0
\]
Sugerencia: calcule primero $\dfrac{\partial r}{\partial x}$, $\dfrac{\partial r}{\partial y}$, $\dfrac{\partial \theta}{\partial x}$, $\dfrac{\partial \theta}{\partial y}$, y use la regla de la cadena para expresar $u_x$, $u_y$, $u_{xx}$, $u_{yy}$ en términos de $u_r$, $u_\theta$, $u_{rr}$, $u_{r\theta}$, $u_{\theta\theta}$.
